\section{Performance analysis}\label{sec:results}
We tested our model on two different datasets: a dataset (Mix-2022) obtained by the traces collected in a controlled environment, as described in Section \ref{sec:setup}, and a dataset (Cui-2020) resulting from a public repository provided by Cui and Umphress \cite{cui2020towards}. The repository is made of traces captured in a Docker platform running a containerized MySQL server as sample application. Random commands were generate to interact with the server. The system call logs were collected under two scenarios: one with regular and clear application execution and the other with the application under attack. Developers implemented the following attacks: Brute Force Login, DockerEscape, Malicious Script, Meterpreter, Remote Shell, SQL Injection, and SQL Misbehavior. The repository contains 20 traces of system call logs for the benign container execution and 10 traces for each tested attack. We refer the readers to the manuscript by Cui and Umphress  \cite{cui2020towards} for further details about their experiment and data collection.


\subsection{Evaluation metrics}
Given a testing dataset consisting of $M_{n}$ samples from the normal behavior class and $M_{a}$  samples from a generic attack class, let $M_{TP}$ be the number of samples from the attack class that were correctly classified as anomalies (true positive), and $M_{FP}$ be the number of samples from the normal behavior class that were incorrectly classified as anomalies (false positive). In order to evaluate the performance of our algorithm, we used the following traditional metrics: (i) the true positive rate $TPR = \frac{M_{TP}}{M_a}$, (ii) the false positive rate $FPR = \frac{M_{FP}}{M_n}$, (iii) the precision $PR = \frac{M_{TP}}{M_{TP}+M_{FP}}$, and (iv) the score $F1=\frac{2\cdot PR\cdot TPR}{PR + TPR}$.


\subsection{Baseline}
We analyze the performance of the proposed solution by compering its results against a traditional models based on (i) Support Vector Machines (SVM) for the multi-class classification in the first stage and (ii) Local Outlier Factor and One-class SVM for second stage outlier detection, receiving as input the bag of system calls computed over a predefined observation period.

\subsubsection{Bag of System Calls}
A ``bag of system calls'' (BoSC) is a representation of the list of system calls observed by the kernel during a fixed period of time. Let  $S=\{s_i\}_i=1,\dots,N$ be the set of the $N$ types of different system calls for the operating system, and $L=(l_1,\dots,l_{K})$ be the sequence of $K$ system calls observed during the period of time $T=[t_{start},t_{end})$, with $l_k\in S$ for $k_1,\dots,K$. The bag of system calls $B_T$ for the time period $T$ is defined as $B_T=(b_1,b_2,\dots,b_N)$ where $b_i$ is the number of occurrences of the system call $s_j$ in the sequence $L$. BoSC-based representations are widely used in intrusion detection problems, as highlighted in Section \ref{sec:relworks}.

\subsubsection{Support Vector Machines}
Support Vector Machine (SVM) is a state-of-the-art supervised machine learning model commonly used as baseline in the intrusion detection domain, since it provides good and stable performance in classification and anomaly detection problems.

In an SVM classifier, samples in a training dataset are represented as point in an n-dimensional space. The SVM training algorithm aims at identifying the hyperplanes that partition the hyperspace in subspaces and separate the classes in the hyperspace. The hyperplanes are defined starting from the support vectors, i.e. data points in the training set that are closer to an hyperplane. By adding or removing support vectors, the algorithm influences the position and orientation of the hyperplanes so as to maximize the distance between the support vectors and the hyperplanes (margin). The hyperplanes resulting from the training algorithm will be the decision boundaries to determine the class of a given data sample.

In One-class SVM outlier detection, the algorithm aims at defining the hyperplanes that separate all the samples in the training set from the origin of the hyperspace.

We refer the readers to the original manuscripts by Boser  \cite{boser1992training} and Scholkopf \cite{scholkopf1999support} for further details about SVM-based classification and outlier detection.

\subsubsection{Local Outlier Factor}

The Local Outlier Factor (LOF) algorithm is an unsupervised anomaly detection method proposed by Breuning \etal \cite{breunig2000lof}. The algorithm computes for each data point an outlier score that measures the relative density deviation between the data point and its local neighbors; when the local density of the data point is lower than the density of its nearest neighbors, the data point is considered an outlier. 
We refer the readers to the original manuscripts by Breuning \etal \cite{breunig2000lof} for further details.

\subsection{Dataset creation and processing}
Starting from the original system call logs recorded and locally  stored, we randomly selected 10 not overlapping slots of duration $T=10$ sec for each trace in our repository so we could generate dataset ``Mix-2022'' made of 10 thousand samples for each class. Instead, for the traces available in Cui-2020, we randomly selected up to 15 not overlapping  slots of duration $T=10$ sec for each trace, obtaining a dataset made of 733 samples. 

Both datasets were split into training and testing datasets. In dataset Mix-2022, the samples belonging to the ``benign'' classes were split according to the following proportions: 70\% into the training and 30\% into the testing dataset; all the attack samples were allocated into the testing dataset. We derived four different training sets: S-set containing 70\% samples of the training set to train the (supervised) classifier, and the U(k)-sets (for $k=1,2,3$)  made of (i) all the samples in the training set belonging to the class $k$ and 2.5\% of the remaining samples in the training set belonging to the other classes (contamination). The U(k)-sets are used to train the unsupervised anomaly detection algorithm.

For dataset Cui-2020, we adopted the following process: the samples belonging to the ``benign'' class were split according to the 70-30 proportion, and all the attack samples were allocated into the testing dataset, as for dataset Mix-2022. The training dataset was contaminated with samples randomly generated and made of features drawn by independent probability mass functions obtained by the feature values in the training set. We set a contamination rate equal to 5\%.


\begin{figure*}\centering
    \subfloat[EoF - Dataset ``Mix-2022.'']{\includegraphics[width=0.48\columnwidth, trim=0cm 0cm 0cm 0cm, clip]{ROC-dataset-1m.pdf}\label{fig:rocgid}}\quad
    \subfloat[SVM - Dataset ``Mix-2022.'']{\includegraphics[width=0.48\columnwidth, trim=0cm 0cm 0cm 0cm, clip]{ROC-dataset-1svm.pdf}\label{fig:rocbasesvm}}\quad
    \subfloat[LOF - Dataset ``Mix-2022.'']{\includegraphics[width=0.48\columnwidth, trim=0cm 0cm 0cm 0cm, clip]{ROC-dataset-1lof.pdf}\label{fig:rocbaselof}}\quad
    \subfloat[Dataset ``Cui-2020.'']{\includegraphics[width=0.48\columnwidth, trim=0cm 0cm 0cm 0cm, clip]{ROC-dataset-3.pdf}\label{fig:roccui}}
    \caption{ROC function for the three models.}\label{fig:roc2022}

\end{figure*}

\begin{table*}
\setlength{\tabcolsep}{6pt}
    \centering
    \caption{Performance comparison of the three models for dataset ``Mix-2022.''}\label{tab:tfprfirst}
    \begin{tabular}{ lccccccccc}
        \toprule
                &  \multicolumn{3}{c}{Ensemble of forests} & \multicolumn{3}{c}{One-class SVM} & \multicolumn{3}{c}{LOF}\\
        \midrule
                &  TPR & PR & F1 &  TPR & PR& F1 & TPR & PR& F1 \\
        \midrule
Data Analytics - C-mining	&	0.881	&	0.975	&	0.926	&	0.235	&	0.910	&	0.374	&	0.150	&	0.960	&	0.260	\\
Media Streaming - C-mining	&	1.000	&	0.982	&	0.991	&	1.000	&	0.982	&	0.991	&	0.996	&	0.818	&	0.054	\\
Web Search - C-mining	&	0.999	&	0.966	&	0.982	&	0.369	&	0.982	&	0.528	&	0.276	&	0.999	&	0.998	\\
Data Analytics - Backdoor	&	0.661	&	0.967	&	0.785	&	0.085	&	0.786	&	0.154	&	0.028	&	0.989	&	0.106	\\
Media Streaming - Backdoor	&	0.678	&	0.974	&	0.800	&	0.159	&	0.897	&	0.270	&	0.056	&	0.970	&	0.430	\\
Web Search - Backdoor	&	0.857	&	0.960	&	0.906	&	0.004	&	0.118	&	0.008	&	0.004	&	0.309	&	0.008	\\
        \bottomrule
    \end{tabular}
\end{table*}

\begin{table*}
\setlength{\tabcolsep}{6pt}
    \centering
    \caption{Performance comparison of the three models for dataset ``Cui-2020.''}\label{tab:tfprsecond}
    \begin{tabular}{ lccccccccc}

        \toprule
                &  \multicolumn{3}{c}{Ensemble of forests} & \multicolumn{3}{c}{One-class SVM} & \multicolumn{3}{c}{LOF}\\
        \midrule
                &  TPR & PR & F1 &  TPR & PR& F1&  TPR & PR& F1 \\
        \midrule
Brute Force Login	&	0.493	&	0.949	&	0.649	&	0.093	&	0.438	&	0.154	&	0.067	&	0.357	&	0.112	\\
DockerEscape	&	0.769	&	0.938	&	0.845	&	0.410	&	0.640	&	0.500	&	0.590	&	0.719	&	0.648	\\
MaliciousScript	&	0.711	&	0.964	&	0.818	&	1.000	&	0.894	&	0.944	&	1.000	&	0.894	&	0.944	\\
Meterpreter	&	1.000	&	0.958	&	0.979	&	1.000	&	0.836	&	0.911	&	1.000	&	0.836	&	0.911	\\
RemoteShell	&	0.973	&	0.973	&	0.973	&	0.959	&	0.886	&	0.921	&	0.959	&	0.886	&	0.921	\\
SQLInjection	&	0.610	&	0.959	&	0.746	&	0.961	&	0.892	&	0.925	&	0.857	&	0.880	&	0.868	\\
SQLMisbehavior	&	0.863	&	0.969	&	0.913	&	0.014	&	0.100	&	0.024	&	1.000	&	0.890	&	0.942	\\
        \bottomrule
    \end{tabular}
\end{table*}



\begin{table}
\setlength{\tabcolsep}{6pt}
    \centering
    \caption{Comparison of the FPRs.}\label{tab:fprates}
    \begin{tabular}{ lccc}
        \toprule
                &  EoF & SVM & LOF\\
        \midrule
Data Analytics - Benign	&	0.045	&	0.046	&	0.012	\\												
Media Streaming - Benign	&	0.036	&	0.037	&	0.001	\\												
Web Search  - Benign	&	0.071	&	0.061	&	0.017	\\												
\midrule																			
Cui-2020	&	0.024	&	0.108	&	0.108	\\												


        \bottomrule
    \end{tabular}
\end{table}


\subsection{Results}
We implemented the models by using sklearn-learn library v0.24.2 in Python v3.5.2. Original Python implementation of anonymous walk embeddings algorithm \cite{ivanov2018anonymous} was adopted to translate the graph representation of system calls into feature vectors. We used anonymous walks of length 4 which produced vectors of 15 features. Both random and isolation forests were initialized with 100 estimators each. The SVM and LOF models were initialized with default configuration.

Tables \ref{tab:tfprfirst} and \ref{tab:tfprsecond} present the results obtained respectively for datasets Mix-2022 and Cui-2020. The tables list the TPR, PR, and F1 computed for each attack separately, while the FPRs are shown in Table \ref{tab:fprates}. The FPRs are computed for each workload in dataset Mix-2022, separately. The performance obtained for the ensemble of random and isolation forests is compared against that obtained by the SVMs -based algorithm.

The results show that the Ensemble of random and isolation Forests (EoF) with graph-based features outperforms the SVMs snf LOF -based algorithm with bag of the system calls as input, for almost all of the attacks in both datasets. This is also confirmed by the ROC curves shown in Figures \ref{fig:roc2022}. The curves, which are obtained by varying the decision threshold in the unsupervised algorithms, highlight the better tradeoff between TPRs and FPRs resulting from the ensemble of forests, since all its curves are located quite close to the upper-left hand corner, while the two baselines produce curves more below with fluctuating trends. By observing the results we can also notice that the detection capability may vary with the type of workload tested; this outcome can be explained by considering that the behavior of some container workloads may hide an underlying attack.

Although the quite high average detection rates, some attacks may yet be able to pass undetected, e.g., true positive rates below 0.7 are resulted against the Backdoor and SQL Injection attacks, and below 0.5 against the Brute Force Login attack.

Finally, there are two attacks (Malicious Script and SQL Injection) where the SVMs-based system outperforms the ensemble of forests in terms of detection rate, nonetheless, the corresponding precision and F1 scores are lower, which is due to the high FPR (Table \ref{tab:fprates}) obtained by the SVMs-based model for dataset Cui-2020.

